
\slajdid{05-s051}
\begin{frame}[label=05-s051]
\gusttekst\setlength{\abovedisplayskip}{4pt}\setlength{\belowdisplayskip}{4pt}
Primer sa četvorobitnim registrima i notacijama indeksa iz brojnih sistema
\[Q_t\{c_0.c_{-1}c_{-2}c_{-3}c_{-4}c_{-5}c_{-6}c_{-7}\}\to c_0.c_{-1}c_{-2}c_{-3}0000\to c_0.c_{-1}c_{-2}c_{-3}\]
bez obzira da li je broj negativan ili pozitivan.\\Uočiti
\[\begin{aligned}c_0.c_{-1}c_{-2}c_{-3}&=c_0.c_{-1}c_{-2}c_{-3}0000\\&=c_0.c_{-1}c_{-2}c_{-3}c_{-4}c_{-5}c_{-6}c_{-7}-0.000c_{-4}c_{-5}c_{-6}c_{-7}\end{aligned}\]
Najveća apsolutna greška je u ovom primeru je $0.0001111=\dfrac{15}{2^7}$\par\vspace{3mm}
U opštem slučaju $\textcolor{DEplava}{c_{n-1}c_{n-2}\ldots c_1c_0=11\ldots11}$\quad$\textcolor{DEplava}{\varepsilon=\dfrac{2^n-1}{2^{2n-1}}}$\par\vspace{3mm}
Odsecanjem se rezultat množenja smanjuje.\\
Uočiti da se po apsolutnoj vrednosti pozitivni brojevi smanjuju\\a negativni po apsolutnoj vrednosti rastu
\end{frame}
